\chapter{Multiplicar números relativos}\label{ch:g8:negprod}

El año pasado sumamos y restamos \omterm{def:g7:negatives:def}{números relativos}
(\cref{ch:g7:negatives}); queda multiplicarlos y dividirlos. Una
pregunta domina el capítulo: ¿por qué «menos por menos» tiene que ser
más? Damos una razón, no solo una regla.

\section{Las reglas de los signos}

\begin{theorem}[Signo de un producto]\label{thm:g8:negprod:rules}
El \omterm{def:g2:mult:def}{producto} de dos \omterm{def:g7:negatives:def}{números relativos} tiene una distancia al \omterm{def:g1:counting:zero}{cero} igual al
\omterm{def:g2:mult:def}{producto} de las distancias, y su signo lo da:
\begin{center}
\renewcommand{\arraystretch}{1.3}
\begin{tabular}{c|cc}
$\times$ & $+$ & $-$ \\
\hline
$+$ & $+$ & $-$ \\
$-$ & $-$ & $+$ \\
\end{tabular}
\end{center}
\emph{Signos iguales: \omterm{def:g2:mult:def}{producto} positivo. Signos opuestos: \omterm{def:g2:mult:def}{producto} negativo.}
La misma tabla rige los \omterm{def:g3:division:remainder}{cocientes}.
\end{theorem}

\begin{proof}[Por qué $(-1) \times (-1) = +1$]
Primero, $3 \times (-4)$ significa $(-4) + (-4) + (-4) = -12$: un positivo
por un negativo es negativo. Ahora observa el patrón:
\[
3 \times (-4) = -12, \quad
2 \times (-4) = -8, \quad
1 \times (-4) = -4, \quad
0 \times (-4) = 0 .
\]
En cada paso, el primer factor baja $1$ y el resultado \emph{sube}
$4$. Continuando el patrón un paso más:
\[
(-1) \times (-4) = 0 + 4 = +4 .
\]
Cualquier otra elección rompería la regularidad de la aritmética
(la distributividad). Así que menos por menos tiene que ser más.
\end{proof}

\begin{omfigure}
\begin{tikzpicture}
\begin{axis}[omaxis,
  width=7.8cm, height=5.4cm,
  xmin=-2.6, xmax=3.6, ymin=-14, ymax=10,
  xlabel={$k$}, ylabel={$k \times (-4)$},
  xtick={-2,-1,1,2,3}, ytick={-12,-8,-4,4,8}]
  \addplot[omDef, thick, domain=-2.3:3.3] {-4*x};
  \addplot[omDef, only marks, mark=*, mark size=1.6pt] coordinates
    {(3,-12) (2,-8) (1,-4) (0,0) (-1,4) (-2,8)};
\end{axis}
\end{tikzpicture}

{\small El argumento del patrón en una imagen: los puntos
$(k,\ k \times (-4))$ están alineados, y prolongar la recta más allá de $k = 0$
obliga a $(-1)\times(-4) = +4$ y $(-2)\times(-4) = +8$.}
\end{omfigure}

\begin{example}\label{ex:g8:negprod:products}
\begin{align*}
(-7) \times (+6) &= -42 && \text{(signos opuestos)} \\
(-5) \times (-8) &= +40 && \text{(signos iguales)} \\
(+2.5) \times (-4) &= -10 \\
(-36) \div (-9) &= +4 \\
(+15) \div (-2) &= -7.5 .
\end{align*}
\end{example}

\section{Productos de varios factores}

\begin{proposition}[Signo de un producto largo]\label{prop:g8:negprod:many}
El signo de un \omterm{def:g2:mult:def}{producto} de varios factores no nulos depende solo del
número de factores negativos:
\begin{itemize}
  \item un número \emph{\omterm{def:g3:numbers:evenodd}{par}} de factores negativos da un \omterm{def:g2:mult:def}{producto}
        positivo;
  \item un número \emph{\omterm{def:g3:numbers:evenodd}{impar}} da un \omterm{def:g2:mult:def}{producto} negativo.
\end{itemize}
\end{proposition}

\begin{proof}
Los factores negativos se emparejan y cada pareja aporta un signo positivo
(\cref{thm:g8:negprod:rules}); un recuento \omterm{def:g3:numbers:evenodd}{impar} deja un factor negativo
sin pareja, que vuelve negativo todo el \omterm{def:g2:mult:def}{producto}.
\end{proof}

\begin{example}\label{ex:g8:negprod:many}
$(-2) \times (+3) \times (-5) \times (-1)$: tres factores negativos
(\omterm{def:g3:numbers:evenodd}{impar}), así que el \omterm{def:g2:mult:def}{producto} es negativo; distancias: $2 \times 3 \times 5
\times 1 = 30$. Resultado: $-30$. Decide \emph{primero} el signo y después
multiplica las distancias --- dos tareas fáciles en vez de una propensa
a errores.
\end{example}

\begin{example}[Potencias de negativos]\label{ex:g8:negprod:powers}
$(-2)^2 = (-2) \times (-2) = +4$ y
$(-2)^3 = (-2)^2 \times (-2) = -8$: los exponentes \omterm{def:g3:numbers:evenodd}{pares} dan valores
positivos y los \omterm{def:g3:numbers:evenodd}{impares} conservan el signo. Cuidado con la notación:
$(-3)^2 = 9$ pero $-3^2 = -(3 \times 3) = -9$.
\end{example}

\section{Calcular con las cuatro operaciones}

\begin{method}[Calcular con seguridad con números relativos]\label{met:g8:negprod:compute}
\begin{enumerate}
  \item respeta la jerarquía del \cref{ch:g7:priorities}: paréntesis,
        después $\times$ y $\div$, y después $+$ y $-$;
  \item en cada multiplicación o división, halla \emph{primero el
        signo} y después la distancia;
  \item reescribe un paso por línea.
\end{enumerate}
\end{method}

\begin{example}\label{ex:g8:negprod:compute}
\begin{align*}
5 - 3 \times (-4)
&= 5 - (-12) && \text{(primero el producto: signos opuestos)}\\
&= 5 + 12 = 17 .
\end{align*}
\begin{align*}
\frac{(-6) \times (-10)}{-4}
&= \frac{+60}{-4} && \text{(numerador: signos iguales)}\\
&= -15 && \text{(cociente: signos opuestos).}
\end{align*}
\end{example}

\begin{example}[Sustituir valores negativos]\label{ex:g8:negprod:substitute}
Evalúa $E = 3x^2 - 2x$ para $x = -5$, con paréntesis alrededor del valor:
\[
E = 3 \times (-5)^2 - 2 \times (-5)
= 3 \times 25 - (-10)
= 75 + 10 = 85 .
\]
Los paréntesis de $(-5)^2$ son imprescindibles: sin ellos, el cuadrado
se aplicaría solo al $5$.
\end{example}

\section{Ejercicios}

\begin{exercise}[$\star$]\label{exo:g8:negprod:1}
Calcula:
\[
(-8) \times (+7), \qquad
(-9) \times (-6), \qquad
(+12) \times (-0.5), \qquad
(-1) \times (-1) \times (-1).
\]
\end{exercise}

\begin{exercise}[$\star$]\label{exo:g8:negprod:2}
Calcula:
\[
(-63) \div (+9), \qquad
(-8) \div (-16), \qquad
\frac{45}{-5}, \qquad
\frac{-7.2}{-8} .
\]
\end{exercise}

\begin{exercise}[$\star$]\label{exo:g8:negprod:3}
Da solamente el \emph{signo} de cada \omterm{def:g2:mult:def}{producto}, sin calcularlo:
\[
(-3) \times (-7) \times (+2) \times (-5); \qquad
(-1)^{10}; \qquad
(-2)^{7}; \qquad
(-4) \times 0 \times (-6).
\]
\end{exercise}

\begin{exercise}[$\star$]\label{exo:g8:negprod:4}
Calcula:
\[
(-2)^4, \qquad
-2^4, \qquad
(-3)^3, \qquad
(-1)^{2026} .
\]
\end{exercise}

\begin{exercise}[$\star$]\label{exo:g8:negprod:5}
Calcula paso a paso, respetando la jerarquía:
\[
7 + 2 \times (-6), \qquad
(-4) \times (5 - 9), \qquad
18 \div (-3) - (-2) .
\]
\end{exercise}

\begin{exercise}[$\star$]\label{exo:g8:negprod:6}
Calcula:
\[
\frac{(-5) \times (+8)}{-4}, \qquad
\frac{(-3) \times (-6)}{(-2) \times (+9)} .
\]
\end{exercise}

\begin{exercise}[$\star$]\label{exo:g8:negprod:7}
Evalúa para $x = -3$:
\[
5x, \qquad x^2, \qquad -x^2, \qquad 2x^2 + 4x, \qquad (2x)^2 .
\]
\end{exercise}

\begin{exercise}[$\star\star$]\label{exo:g8:negprod:8}
Completa cada igualdad:
\[
(-6) \times \;?\; = 42, \qquad
\;?\; \div (-4) = -2.5, \qquad
(-5) \times \;?\; = -1 .
\]
\end{exercise}

\begin{exercise}[$\star\star$]\label{exo:g8:negprod:9}
¿Verdadero o falso? Justifícalo o da un contraejemplo.
\begin{enumerate}
  \item el \omterm{def:g2:mult:def}{producto} de dos \omterm{def:g7:negatives:def}{números relativos} es siempre al menos tan
        grande como su \omterm{def:g1:addition:def}{suma};
  \item el cuadrado de un \omterm{def:g7:negatives:def}{número relativo} nunca es negativo;
  \item si un \omterm{def:g2:mult:def}{producto} de tres factores es positivo, los tres factores
        son positivos.
\end{enumerate}
\end{exercise}

\begin{exercise}[$\star\star$]\label{exo:g8:negprod:10}
Cada respuesta incorrecta de un concurso cuenta $-3$ puntos y cada acierto,
$+5$. Zoe respondió a $20$ preguntas y sacó $36$ puntos. ¿Cuántas
respuestas fueron correctas? (Prueba valores o plantea el cálculo
$5r - 3(20 - r) = 36$.)
\end{exercise}

\begin{exercise}[$\star\star\star$]\label{exo:g8:negprod:11}
Usando la distributividad (como en la demostración del
\cref{thm:g8:negprod:rules}), desarrolla y justifica cada paso de
\[
0 = (-1) \times 0 = (-1) \times \bigl(1 + (-1)\bigr)
= (-1) \times 1 + (-1) \times (-1),
\]
y concluye que $(-1) \times (-1)$ tiene que valer $+1$.
\end{exercise}

\section{Problema: Por qué menos por menos es más}

\begin{problem}[{Problema de fin de semana --- las reglas de los signos deducidas de
la distributividad y la suma alternada $1 - 2 + 3 - 4 + \dots$}]
\label{pb:g8:negprod:1}
La demostración del \cref{thm:g8:negprod:rules} continuaba un patrón y
afirmaba que «cualquier otra elección rompería la distributividad». Este
problema hace exacta esa afirmación: partiendo solo de la distributividad
(\cref{thm:g7:priorities:distributivity}), \emph{demostrarás} las
reglas de los signos --- sin dibujos ni patrones, como lo hace el álgebra ---
y después las pondrás a trabajar con las potencias de $-1$ y con una \omterm{def:g1:addition:def}{suma} famosa de
signos alternados. En todo momento, $a$ y $b$ son \omterm{def:g7:negatives:def}{números relativos}; damos
por sabidas las reglas de la \omterm{def:g1:addition:def}{suma} del \cref{ch:g7:negatives}, que
$1 \times a = a$ y que un \omterm{def:g2:mult:def}{producto} se puede calcular en cualquier orden
(así que la distributividad se aplica por los dos lados de un \omterm{def:g2:mult:def}{producto}).

\medskip
\noindent\textbf{Parte I --- Las reglas de los signos son teoremas.}
\begin{enumerate}
  \item demuestra que $0 \times a = 0$ para todo \omterm{def:g7:negatives:def}{número relativo} $a$.
        (Escribe $0 = 0 + 0$, desarrolla $(0 + 0) \times a$ y pregúntate qué
        número, sumado a $0 \times a$, devuelve $0 \times a$.)
  \item desarrolla $\bigl(1 + (-1)\bigr) \times a$ para mostrar que
        $(-1) \times a + a = 0$, y concluye que
        \[
        (-1) \times a = -a :
        \]
        multiplicar por $-1$ da el opuesto. Compáralo con el
        \cref{exo:g8:negprod:11}, que es el caso $a = -1$;
  \item deduce que $(-a) \times b = -(a \times b)$: un signo menos
        sale intacto de un \omterm{def:g2:mult:def}{producto};
  \item deduce que $(-a) \times (-b) = a \times b$: dos signos menos
        se cancelan;
  \item todo \omterm{def:g7:negatives:def}{número relativo} es su distancia al \omterm{def:g1:counting:zero}{cero} con un signo
        delante: $a = +d$ o $a = -d$. Explica cómo las preguntas 3 y~4
        demuestran los cuatro casos de la tabla de signos del
        \cref{thm:g8:negprod:rules}, distancias incluidas.
\end{enumerate}

\medskip
\noindent\textbf{Parte II --- Las potencias de $-1$.}
\begin{enumerate}[resume]
  \item calcula $(-1)^2$, $(-1)^3$, $(-1)^4$ y $(-1)^5$, y da después,
        justificándolo con la \cref{prop:g8:negprod:many}, el
        valor de $(-1)^n$ para todo número entero $n \geq 1$;
  \item da (sin calcular ninguna distancia) el signo de
        \[
        (-1) \times (-2) \times (-3) \times \dots \times (-10),
        \]
        y escribe la distancia de ese \omterm{def:g2:mult:def}{producto} como un \omterm{def:g2:mult:def}{producto} de números
        enteros, sin desarrollarlo;
  \item demuestra que $(-a)^2 = a^2$ y que $(-a)^3 = -a^3$. ¿Para qué
        exponentes sobrevive el signo menos?
  \item usa las reglas de los signos para mostrar que el cuadrado de un \omterm{def:g7:negatives:def}{número
        relativo} nunca es negativo, y deduce que ningún \omterm{def:g7:negatives:def}{número relativo}
        $x$ cumple $x^2 = -4$;
  \item Zoe afirma: «$(-1)^{2026} + (-1)^{2027} = 0$ y, más en
        general, dos potencias consecutivas de $-1$ siempre se cancelan». ¿Tiene
        razón? Justifícalo.
\end{enumerate}

\medskip
\noindent\textbf{Parte III --- La \omterm{def:g1:addition:def}{suma} alternada.}
Con las reglas de los signos aseguradas, podemos calcular \omterm{def:g1:addition:def}{sumas} que mezclan los dos
signos. Para un número entero $n \geq 1$, sea
\[
A_n = 1 - 2 + 3 - 4 + \dots
\]
la \omterm{def:g1:addition:def}{suma} de los números enteros de $1$ a $n$ con signos
alternados: el último término es $+n$ si $n$ es \omterm{def:g3:numbers:evenodd}{impar} y $-n$ si $n$ es
\omterm{def:g3:numbers:evenodd}{par}. Así, $A_1 = 1$ y $A_2 = 1 - 2 = -1$.
\begin{enumerate}[resume]
  \item calcula $A_3$, $A_4$, $A_5$, $A_6$ y $A_7$. ¿Qué
        conjeturas?
  \item supón que $n$ es \omterm{def:g3:numbers:evenodd}{par}. Agrupa los términos de dos en dos,
        $(1 - 2) + (3 - 4) + \dots$, cuenta las parejas y demuestra que
        $A_n = -\frac{n}{2}$.
  \item supón que $n$ es \omterm{def:g3:numbers:evenodd}{impar}. Explica por qué $A_n = A_{n-1} + n$ y
        deduce de la pregunta anterior que
        $A_n = \frac{n + 1}{2}$.
  \item calcula $A_{100}$, $A_{101}$ y $A_{2026}$;
  \item Alma calculó $A_{101} - A_{100} = 51 - (-50) = 101$ y le
        sorprendió obtener exactamente $101$. ¿Debería haberle sorprendido?
        Explícalo en una frase y di el valor de
        $A_{2027} - A_{2026}$ sin ningún cálculo más.
\end{enumerate}
\end{problem}
